
\documentclass{article}
\usepackage{colortbl}
\usepackage{makecell}
\usepackage{multirow}
\usepackage{supertabular}

\begin{document}

\newcounter{utterance}

\centering \large Interaction Transcript for game `imagegame', experiment `random\_grids', episode 17 with reflexion{-}llama{-}3.1{-}8b{-}instruct.
\vspace{24pt}

{ \footnotesize  \setcounter{utterance}{1}
\setlength{\tabcolsep}{0pt}
\begin{supertabular}{c@{$\;$}|p{.15\linewidth}@{}p{.15\linewidth}p{.15\linewidth}p{.15\linewidth}p{.15\linewidth}p{.15\linewidth}}
    \# & $\;$A & \multicolumn{4}{c}{Game Master} & $\;\:$B\\
    \hline

    \theutterance \stepcounter{utterance}  
    & & \multicolumn{4}{p{0.6\linewidth}}{
        \cellcolor[rgb]{0.9,0.9,0.9}{
            \makecell[{{p{\linewidth}}}]{
                \texttt{\tiny{[P1$\langle$GM]}}
                \texttt{You have attempted this same task before and did not succeed. These are the plans you wrote after those attempts.} \\
\texttt{Trial \#1: New plan: I got stuck in an infinite loop of putting Rs in the first row because I didn't consider the existing Rs in the grid. Next time, I will make sure to check the grid for existing Rs before placing a new one. I will also pay attention to the format of the command, which requires a specific description of the action, such as "Put an R in first row fifth column". To avoid violating the rules, I will make sure to provide a complete and accurate description of the action, including the row and column numbers.} \\
\\ 
\texttt{Trial \#2: New plan: I got stuck in an infinite loop of putting Rs in the first row because I didn't consider the existing Rs in the grid. However, I also failed to follow the format of the command, which requires a specific description of the action, such as "Put an R in first row fifth column". To avoid violating the rules, I will make sure to provide a complete and accurate description of the action, including the row and column numbers. Next time, I will check the grid for existing Rs and pay attention to the format of the command. I will also consider the size of the grid and the existing letters when planning my next move. For example, I will start by placing an R in a row where there is no existing R, and then I will move to the next step.} \\
\\ 
\\ 
\texttt{Let us do the following task. The goal is to fill an empty grid that looks like this:} \\
\\ 
\texttt{▢ ▢ ▢ ▢ ▢} \\
\texttt{▢ ▢ ▢ ▢ ▢} \\
\texttt{▢ ▢ ▢ ▢ ▢} \\
\texttt{▢ ▢ ▢ ▢ ▢} \\
\texttt{▢ ▢ ▢ ▢ ▢} \\
\\ 
\texttt{A filled grid below is 5 by 5 and can look like this:} \\
\\ 
\texttt{▢ ▢ ▢ ▢ ▢} \\
\texttt{▢ ▢ E ▢ ▢} \\
\texttt{▢ ▢ ▢ ▢ ▢} \\
\texttt{▢ ▢ ▢ ▢ ▢} \\
\texttt{X X X X X} \\
\\ 
\texttt{I want you to describe this grid to me, step by step. You don't need to describe the empty squares, which are denoted with "▢". Only describe the location of letters in the grid. Then you wait for me to say "What is your next command?", and then you continue with the next step. Take the size of the grid into consideration while giving commands. You can give only one command at each turn but it can refer to single or multiple cells, rows, columns or even the whole grid at once. When you have described everything, you say "Command: DONE". Only give the command and omit everything else.} \\
\\ 
\texttt{For the filled grid above, here are the example steps.} \\
\\ 
\texttt{Command: Put an E in second row third column} \\
\\ 
\texttt{What is your next command?} \\
\texttt{Command: Fill the last row with X} \\
\\ 
\texttt{What is your next command?} \\
\texttt{Command: DONE} \\
\\ 
\texttt{Ok. Please do this for the following example, which is a 5 by 5 grid.} \\
\texttt{R ▢ ▢ R ▢} \\
\texttt{▢ ▢ ▢ ▢ ▢} \\
\texttt{▢ ▢ ▢ R ▢} \\
\texttt{▢ ▢ ▢ R ▢} \\
\texttt{▢ ▢ ▢ ▢ R} \\
\texttt{What is your next command?} \\
            }
        }
    }
    & \\ \\

    \theutterance \stepcounter{utterance}  
    & \multicolumn{4}{p{0.6\linewidth}}{
        \cellcolor[rgb]{0.9,0.9,0.9}{
            \makecell[{{p{\linewidth}}}]{
                \texttt{\tiny{[P1$\rangle$GM]}}
                \texttt{There is an R in the first row first column and fifth column. There is an R in the third row fifth column and fourth row fifth column. There is an R in the fourth row third column and fourth row fourth column. There is an R in the fifth row fifth column.} \\
            }
        }
    }
    & & \\ \\

    \theutterance \stepcounter{utterance}  
    & & & \multicolumn{2}{p{0.3\linewidth}}{
        \cellcolor[rgb]{0.9,0.9,0.9}{
            \makecell[{{p{\linewidth}}}]{
                \texttt{\tiny{[GM$|$GM]}}
                \texttt{Invalid instruction format} \\
            }
        }
    }
    & & \\ \\

    \theutterance \stepcounter{utterance}  
    & \multicolumn{4}{p{0.6\linewidth}}{
        \cellcolor[rgb]{0.95,0.95,0.95}{
            \makecell[{{p{\linewidth}}}]{
                \texttt{\tiny{[P2$\rightarrow$P1]}}
                \texttt{Post{-}episode Reflexion (not sent to GM):} \\
\texttt{New plan: I got stuck in an infinite loop of putting Rs in the first row because I didn't consider the existing Rs in the grid. However, I also failed to follow the format of the command, which requires a specific description of the action, such as "Put an R in first row fifth column". To avoid violating the rules, I will make sure to provide a complete and accurate description of the action, including the row and column numbers. Next time, I will check the grid for existing Rs and pay attention to the format of the command. I will also consider the size of the grid and the existing letters when planning my next move. For example, I will start by placing an R in a row where there is no existing R, and then I will move to the next step. To make progress, I will focus on filling in the rows that have the fewest existing Rs, and I will try to place Rs in columns where there are no existing Rs.} \\
            }
        }
    }
    & & \\ \\

\end{supertabular}
}

\end{document}
